\documentclass[12pt,a4paper]{report}

% ---------- Packages ----------
\usepackage[a4paper,margin=1in]{geometry}
\usepackage[T1]{fontenc}
\usepackage{graphicx}
\usepackage{float}
\usepackage{setspace}
\usepackage{booktabs}
\usepackage{array}
\usepackage{longtable}
\usepackage{hyperref}
\usepackage{pdfpages}
\usepackage{caption}
\usepackage{xcolor}
\usepackage{enumitem}
\usepackage{titlesec}

\graphicspath{{documentation/}}
\setstretch{1.15}
\setlength{\parskip}{3pt}
\setlength{\parindent}{15pt}
\setlist{topsep=3pt,itemsep=2pt,parsep=0pt,partopsep=0pt}
\titlespacing*{\chapter}{0pt}{-10pt}{14pt}
\titlespacing*{\section}{0pt}{10pt}{4pt}
\titlespacing*{\subsection}{0pt}{8pt}{3pt}
\hypersetup{
    colorlinks=true,
    linkcolor=blue,
    urlcolor=blue,
    citecolor=blue
}
\setcounter{secnumdepth}{3}
\setcounter{tocdepth}{2}

% Helper: a diagram gets its own clean page with its caption
\newcommand{\projectdiagram}[3]{%
    \clearpage
    \section{#1}
    #2
    \begin{center}
        \includegraphics[width=0.97\textwidth,height=0.80\textheight,keepaspectratio]{#3}
        \captionof{figure}{#1}
    \end{center}
    \clearpage
}

\begin{document}

% ---------- Title Page ----------
\begin{titlepage}
    \centering
    \vspace*{1cm}
    {\Large\bfseries THAPAR INSTITUTE OF ENGINEERING \& TECHNOLOGY\par}
    \vspace{0.4cm}
    {\large Department of Computer Science and Engineering\par}
    \vspace{1.8cm}
    {\Huge\bfseries THAPAR TRADE\par}
    \vspace{0.5cm}
    {\Large Final Software Engineering Project Report\par}
    \vspace{1.8cm}
    \textbf{Submitted by}\\[0.5cm]
    Rajbir Singh -- 1024030453\\
    Chitesh Jindal -- 1024030454\\
    Ishan Garg -- 1024031095\\[1.3cm]
    \textbf{Submitted to:} Jeelani Asif\\[0.5cm]
    \textbf{Course:} UCS503P -- Software Engineering Lab\\[0.3cm]
    \textbf{Academic Year:} 2026
    \vfill
    Thapar Institute of Engineering \& Technology\\
    Patiala, Punjab
\end{titlepage}

\pagenumbering{roman}

% ---------- Abstract ----------
\chapter*{Abstract}
\addcontentsline{toc}{chapter}{Abstract}
Thapar Trade is a campus marketplace prototype designed for the Thapar community.
It provides a single interface for discovering campus listings, publishing items
for sale, saving favourites, managing a cart, completing checkout through a
simulated payment flow, viewing orders, and tracking order status. The project
also includes listing reports and administrative moderation features.

The frontend is designed to operate as a standalone academic prototype, using
browser local storage for persistence. An optional FastAPI and SQLite backend is
also included for local full-stack demonstration. The report describes the
problem, requirements, system design, implementation, testing, project schedule,
results, limitations, and future scope.

\tableofcontents
\listoffigures
\clearpage
\pagenumbering{arabic}

% ---------- 1 Introduction ----------
\chapter{Introduction}
Thapar Trade is a web-based campus marketplace intended for members of the
Thapar community. The platform brings common peer-to-peer marketplace activities
into one place. Users can browse available items, search and filter listings,
publish and manage their own listings, save items, manage a cart, and view orders.
The project focuses on making campus item discovery and exchange more organised
than relying only on informal messages or separate communication groups.

\section{Project Objectives}
\begin{itemize}
    \item Provide a central marketplace interface for the Thapar community.
    \item Allow users to create, edit, and delete item listings.
    \item Support marketplace browsing, keyword search, and category filtering.
    \item Provide favourites and cart functionality.
    \item Support checkout using a simulated payment flow.
    \item Provide order history and order-status tracking.
    \item Support listing reports and administrative moderation.
\end{itemize}

% ---------- 2 Background ----------
\chapter{Background}
Campus communities regularly exchange items among students. When listings are
shared through informal channels, users may need to search through unrelated
messages, contact sellers separately, and keep track of item details manually.
A dedicated marketplace can organise listings and provide a consistent interface
for browsing and managing exchanges.

Thapar Trade addresses this context with a campus-focused marketplace prototype.
It groups listing discovery, seller information, saved listings, cart operations,
and order-related workflows in one application.

% ---------- 3 Problem Statement ----------
\chapter{Problem Statement}
Students need a convenient way to discover and publish items within their campus
community. Informal channels do not necessarily provide structured listing
details, category-based discovery, saved items, cart management, or a consistent
way to review order status.

The problem addressed by Thapar Trade is the absence of a single organised
campus marketplace interface that brings these activities together for the
Thapar community.

% ---------- 4 Requirements ----------
\chapter{Requirements Analysis}
\section{Functional Requirements}
\begin{enumerate}
    \item Users can access the registration and login interface.
    \item Users can browse marketplace listings and view item details.
    \item Users can search listings by keyword and filter by category.
    \item Users can create, edit, and delete their own listings.
    \item Users can add and remove listings from favourites.
    \item Users can add eligible items to a cart and manage quantities.
    \item Users can proceed through checkout using the simulated payment flow.
    \item Users can view order history and order status.
    \item Users can submit ratings, reviews, and listing reports.
    \item Administrators can access moderation and order-management controls.
\end{enumerate}

\section{Non-Functional Requirements}
\begin{itemize}
    \item The interface should be usable on desktop and mobile screen sizes.
    \item Common marketplace actions should be presented through a consistent UI.
    \item The standalone prototype should retain supported data across browser refreshes.
    \item The codebase should support local development and production build generation.
\end{itemize}

% ---------- 5 Proposed Solution ----------
\chapter{Proposed Solution}
Thapar Trade provides a central web interface for campus marketplace activities.
A user can browse listings, search for an item, apply category filters, save
favourites, and manage eligible products in a cart. A seller can create and
manage listings. The application also provides checkout, simulated payment,
order history, and order-status workflows. Reporting and administrative
moderation features support basic marketplace oversight.

\section{Main Modules}
\begin{description}
    \item[Authentication interface:] Provides the user registration and login flow.
    \item[Marketplace:] Displays listings and supports search and category filtering.
    \item[Listing management:] Supports creating, editing, and deleting listings.
    \item[Favourites and cart:] Supports saving listings and managing eligible cart items.
    \item[Checkout and orders:] Supports checkout with simulated payment and order records.
    \item[Reviews and reports:] Supports feedback and listing-report workflows.
    \item[Administration:] Provides listing moderation and order-management controls.
\end{description}

% ---------- 6 Architecture ----------
\chapter{System Architecture and Technology}
\section{Technology Stack}
\begin{itemize}
    \item \textbf{Frontend:} React.js
    \item \textbf{Programming language:} JavaScript
    \item \textbf{Styling:} CSS
    \item \textbf{Standalone persistence:} Browser local storage
    \item \textbf{Optional backend:} FastAPI
    \item \textbf{Optional database:} SQLite
\end{itemize}

\section{Architecture Overview}
The standalone frontend provides the user interface and marketplace workflows.
Browser local storage is used for persistence in the academic prototype. An
optional FastAPI and SQLite backend is included separately for local full-stack
demonstration. The GitHub Pages frontend build does not require that optional
backend.

\section{Important Prototype Scope}
The frontend checkout uses simulated payment; it does not process a real
financial transaction. Browser local storage is used for standalone persistence.
The prototype authentication and administration mechanisms should not be
interpreted as production-grade identity or access security.

% ---------- 7 Diagrams ----------
\chapter{System Design Diagrams}

\section{Use Case Diagram}
The use case diagram presents the interactions between marketplace users and the main system functions.
\begin{center}
    \includegraphics[width=0.97\textwidth,height=0.80\textheight,keepaspectratio]{User_case_diagram.png}
    \captionof{figure}{Use Case Diagram}
\end{center}
\clearpage

\projectdiagram{Class Diagram}
{The class diagram represents the main classes and their relationships in the system design.}
{Class-diagram.png}

\projectdiagram{Activity Diagram}
{The activity diagram illustrates the flow of activities within the application.}
{Activity_diagram.jpg}

\projectdiagram{Sequence Diagram}
{The sequence diagram shows the sequence of interactions between system participants and components.}
{Sequence-diagram.png}

\projectdiagram{Collaboration Diagram}
{The collaboration diagram presents the participating objects and their interaction relationships.}
{Collaboration-diagram.png}

\projectdiagram{Deployment Diagram}
{The deployment diagram illustrates the deployment arrangement represented in the supplied project design.}
{Deployment-Diagram.png}

\projectdiagram{Entity Relationship Diagram}
{The entity relationship diagram represents the data entities and relationships shown in the supplied design.}
{ER_diagram.png}

\projectdiagram{Data Flow Diagram -- Level 0}
{The Level 0 data flow diagram presents the system boundary and its external interactions.}
{DFD_level0.png}

\projectdiagram{Data Flow Diagram -- Level 1}
{The Level 1 data flow diagram expands the system processes and data flows.}
{DFD_level1.png}

% ---------- 8 Implementation ----------
\chapter{Implementation Overview}
\section{Marketplace and Listings}
The marketplace interface supports listing discovery, keyword search, and
category filtering. Listing management supports adding, editing, and deleting
items. Item details and seller information are presented as part of the
marketplace workflow.

\section{Favourites and Cart}
Users can save listings to favourites and manage eligible items in a cart.
Cart quantity management and restrictions are part of the purchase workflow.

\section{Checkout and Order Management}
The checkout workflow collects order information and creates an order record
using the prototype's simulated payment process. Users can view order history
and order status. No live payment gateway is represented by this prototype.

\section{Reviews, Reports, and Administration}
The application includes ratings and reviews, listing reports, and an
administrative moderation dashboard. Administrative features support review
of listings and order-related status workflows.

% ---------- 9 Testing ----------
\chapter{Testing and Validation}
The project team tested eight major test scenarios covering the principal
application workflows. The detailed test cases, expected results, actual
results, and execution status are included in the attached test-case document.

\clearpage
% Keep Test_Cases.pdf in the same Overleaf folder as main.tex.
\includepdf[pages=-,pagecommand={},fitpaper=true]{documentation/test_cases/Test_Cases.pdf}
\clearpage

% ---------- 10 Gantt ----------
\chapter{Project Schedule and Gantt Chart}
The project schedule records the planned activities, assigned team members,
progress, and dates across the project period. The complete Gantt chart is
included below in its original spreadsheet-template layout.

\clearpage
% Keep Gantt_Chart.pdf in the same Overleaf folder as main.tex.
\includepdf[pages=-,pagecommand={},fitpaper=true]{documentation/gantt_chart/Gantt_Chart.pdf}
\clearpage

% ---------- 11 Results ----------
\chapter{Results}
The delivered prototype brings the main campus marketplace workflows into one
application. Its implemented scope includes marketplace browsing, listing
management, search and filtering, favourites, cart operations, simulated
checkout, order history and tracking, reviews, reports, and administrative
moderation.

The test-case document records the eight major scenarios tested by the project
team. The project also includes the supplied system-design diagrams and the
project schedule.

% ---------- 12 Limitations ----------
\chapter{Limitations and Future Scope}
\section{Limitations}
\begin{itemize}
    \item The frontend checkout uses simulated payment rather than a live payment gateway.
    \item The standalone frontend uses browser local storage for persistence.
    \item The academic prototype's authentication and administration features
    are not production-grade security mechanisms.
    \item The optional backend is provided for local full-stack demonstration;
    the GitHub Pages frontend build is standalone.
\end{itemize}

\section{Future Scope}
\begin{itemize}
    \item Integrate a production backend and hosted database.
    \item Implement secure institutional identity verification and server-side authorization.
    \item Integrate a real payment gateway after completing the required onboarding.
    \item Add reliable notification and delivery-status integrations.
    \item Improve monitoring, audit logging, and deployment security.
\end{itemize}

% ---------- 13 Conclusion ----------
\chapter{Conclusion}
Thapar Trade demonstrates a campus-focused marketplace prototype that combines
item discovery, listing management, favourites, cart operations, checkout,
order tracking, reviews, reports, and administrative moderation. The project
provides a structured digital interface for marketplace activities within the
Thapar community. The included diagrams, test cases, and project schedule
document the design and development work.

% ---------- References ----------
\begin{thebibliography}{9}
\bibitem{react}
React Documentation. \url{https://react.dev/}

\bibitem{fastapi}
FastAPI Documentation. \url{https://fastapi.tiangolo.com/}

\bibitem{sqlite}
SQLite Documentation. \url{https://www.sqlite.org/docs.html}

\bibitem{githubpages}
GitHub Pages Documentation. \url{https://docs.github.com/en/pages}
\end{thebibliography}

\end{document}
